% BEGIN LEAD ARC-0007-MULTINORM 2026-10-08
\section{Finite-state phase potentials for existential ARC closure}
The fixed-return theorem above multiplies a prescribed finite word of
nonnegative affine maps.  A different sufficient certificate permits the
next word to depend on the current phase and achieved pair, while retaining a
finite, exact potential proof.

Let \(G=(Q,E)\) be a fixed finite directed graph.  An edge
\(e:q\to r\) carries \(M_e\ge0\) and \(d_e\ge0\).  Suppose there are
positive row vectors \(v_q\), constants \(0<\lambda<1\), \(c\ge0\), and
\(\theta>0\) such that for every edge
\[
 v_r^{\mathsf T}M_e\le\lambda v_q^{\mathsf T},
 \qquad
 v_r^{\mathsf T}d_e\le c,
 \qquad
 \{z\ge0:v_q^{\mathsf T}z<\theta\}\subseteq\mathcal B,
\]
and \(c/(1-\lambda)<\theta\).  The SAT-specific premise is that for every
phase \(q\) and every \(z\in\mathcal R\setminus\mathcal B\), some outgoing
edge \(e:q\to r\) has an achievable successor
\(z'\in\mathcal R\) with \(z'\le M_ez+d_e\).

\begin{theorem}[PHASE-MULTINORM-ARC1]
Under the displayed hypotheses, any finite achievable starting pair forces
\(\mathcal R\cap\mathcal B\ne\varnothing\).
\end{theorem}
\begin{proof}
Choose one permitted successor while outside \(\mathcal B\).  With the
current phase \(q_j\), put \(V_j=v_{q_j}^{\mathsf T}z_j\).  Positivity and the
edge inequalities give
\[
 V_{j+1}\le\lambda V_j+c
 \le \lambda^{j+1}V_0+
       c\frac{1-\lambda^{j+1}}{1-\lambda}.
\]
The right side is eventually below \(\theta\), so a finite achieved pair
lies in \(\mathcal B\).
\end{proof}

This is an exact envelope theorem and a finite linear-inequality certificate,
not a SAT construction.  Philippe et al.'s multinorm assigns one norm to
each automaton node and characterizes stability of constrained switching
systems \cite[Definition~2.2, Proposition~2.3, Theorem~2.4]{arc0007PhilippeMultinorm}.
That is an exact transfer for the finite graph/potential calculation.
Path-complete Lyapunov functions similarly use graph-indexed inequalities
\cite{arc0007AhmadiPathComplete}.  Neither source supplies the displayed
SAT-successor premise, joint machine realization, or Williams-model endpoint.

The exact rational model uses phases \(A,B\), weights
\[
 v_A=(1,1),\qquad v_B=(1/3,3),\qquad \lambda=2/3,\qquad c=1/100.
\]
Its cross edges have matrices
\(\operatorname{diag}(2,1/5)\) and
\(\operatorname{diag}(1/5,2)\); each is individually unstable, so no common
positive linear potential can contract every edge.  The phase weights
nevertheless satisfy all edge inequalities, including two contracting
self-loops.  The limiting potential is \(3/100<1/3\), and
\(v_q^{\mathsf T}z<1/3\) implies \(t+s<1<2\cos(\pi/7)\).
Exact enumeration checked 1,331,200 accepted trajectory steps through depth
ten, together with illegal-phase, negative-matrix, phase-erasure, oracle,
and fixed-graph controls.

The quantifiers remain existence-level.  Under \(P=NP\), first fix one finite
achievable pair, then select finitely many successor witnesses.  The graph,
weights, matrices, and resulting number of theorem applications are constants;
no universal source-code compiler is required.  A graph, phase set, or
iteration depth growing with input length is a different uniform-algorithm
claim and is not covered.  The SAT-specific successor premise is
\textsc{Open}; therefore ARC closure and P versus NP remain open.
See \path{research/agents/lead/2026-10-08_063447_LEAD_ARC-0007.tex}.
% END LEAD ARC-0007-MULTINORM 2026-10-08
